\chapter{نتایج و تفسیر}
\label{chp:results}
\section{مقدمه}
در این فصل، نتایج اصلی تحقیق و یافته‌های کلیدی بررسی و تحلیل می‌شوند. هدف از این تحقیق، ارزیابی عملکرد مدل‌های یادگیری ماشین در پیش‌بینی حرکات بازارهای مالی با استفاده از داده‌های توییت و برچسب‌گذاری سه‌مانع بود. در این راستا، مدل‌های مختلفی مورد بررسی قرار گرفتند و به تحلیل نتایج به دست آمده از این مدل‌ها پرداخته شد. در این فصل، تأثیرات این مدل‌ها بر دقت پیش‌بینی و استراتژی‌های تجاری تحلیل شده و نقاط قوت و ضعف آن‌ها مورد بررسی قرار می‌گیرد. با توجه به تحلیل‌های انجام شده، این فصل به ارائه نتیجه‌گیری‌های کلیدی و پیشنهادات برای تحقیقات آینده می‌پردازد. نتایج حاصل از این تحقیق می‌تواند به بهبود ابزارهای پیش‌بینی و استراتژی‌های معاملاتی کمک کرده و به درک عمیق‌تری از تأثیر اطلاعات متنی بر بازارهای مالی منجر شود. همچنین، این فصل به بررسی محدودیت‌ها و جهت‌گیری‌های آینده تحقیق پرداخته و راهکارهایی برای ادامه تحقیقات و بهره‌برداری از نتایج این تحقیق ارائه می‌دهد.

\section{نتایج یادگیری و تنظیم دقیق}
در این بخش، به بررسی و ارزیابی مدل‌های مختلف پرداخته می‌شود. ابتدا، تأثیر حذف هر یک از سه اطلاعات متنی اضافه شده به توییت‌ها بر عملکرد مدل آگاه به زمینه مورد بررسی قرار می‌گیرد. سپس، نتایج انسداد لایه‌های مدل و تأثیر آن بر منحنی یادگیری مدل‌ها تحلیل می‌شود. در نهایت،، دلایل استفاده از دو دوره آموزشی برای مدل‌ها و تأثیر بیش‌برازش بر عملکرد مدل‌ها مورد بحث قرار می‌گیرد. این تحلیل‌ها به منظور بهبود عملکرد مدل‌ها و جلوگیری از مشکلات احتمالی مانند بیش‌برازش انجام شده‌اند.

\subsection{بیش‌برازش}
در این بخش، دلایل استفاده از دو دوره\LTRfootnote{Epoch} آموزشی برای مدل‌ها توضیح داده شده است. یکی از دلایل این انتخاب، پیروی از روش‌های پیشنهادی موجود در این حوزه است. دلیل دیگر این است که همان‌طور که در شکل \ref{fig:overfit} مشاهده می‌شود، پس از گذشت مدت کوتاهی از دوره دوم، مدل به نتایج نهایی خود می‌رسد و نسبتاً ثابت می‌ماند. ادامه آموزش پس از این نقطه منجر به بیش‌برازش مدل می‌شود. همان‌طور که مشاهده می‌شود، هر سه مدل تنظیم‌ شده پس از دوره دوم به امتیاز نهایی خود می‌رسند و ثابت می‌مانند. بیش‌برازش در مدل‌های زبانی بزرگ زمانی رخ می‌دهد که مدل الگوهای موجود در داده‌های آموزشی را بیش از حد یاد می‌گیرد، از جمله نویز و داده‌های غیرمعمول که به داده‌های جدید تعمیم نمی‌یابند. بهترین روش‌ها برای جلوگیری از بیش‌برازش شامل استفاده از تکنیک‌هایی مانند اعتبار‌سنجی متقابل، روش‌های  مانند حذف تصادفی\LTRfootnote{Dropout}، توقف زودهنگام\LTRfootnote{Early Stopage}زمانی که عملکرد روی مجموعه اعتبارسنجی شروع به کاهش می‌کند (همان‌طور که در این نمودار نشان داده شده است)، و استفاده از مدل‌های ساده‌تر در صورت امکان است \cite{ozdemir2023quick}.
\begin{figure}
    \centering
    \includegraphics[width=1\linewidth]{figures/f1_overfit.png}
    \caption{بیش‌برازش مدل‌های یادگیری بعد از دوره دوم یادگیری}
    \label{fig:overfit}
\end{figure}

\subsection{نتایج حذف ویژگی}
جدول \ref{tbl:ablation_study} نتایج مطالعه حذف ویژگی‌ها\LTRfootnote{Ablation Study} را نشان می‌دهد که در آن هر یک از سه اطلاعات متنی اضافه شده به توییت‌ها حذف شده‌اند تا تأثیر آن‌ها بر عملکرد یادگیری در مدل آگاه به زمینه بررسی شود. همان‌طور که مشاهده می‌شود، حذف برچسب پیشین منجر به کاهش قابل توجهی در همه معیارهای ارزیابی شده است، به طوری که صحت به \lr{43.65\%} کاهش یافته است. این نشان می‌دهد که برچسب پیشین نقش بسیار مهمی در پیش‌بینی مدل دارد. این در حالیست که حذف کردن دیگر موارد تقریبا نتایج یکسانی را با مدل نهایی به دست می‌آورند. این نتایج نشان می‌دهد که استفاده از برچسب پیشین به دلیل استفاده از برچسب‌گذاری سه‌مانع که شامل پنجره‌های برچسب‌ها است، تأثیر زیادی بر دقت پیش‌بینی دارد. لازم به ذکر است که منظور از برچسب پیشین، برچسب پنجره پیشین است و مربوط به روز قبل نیست (برچسب روز قبل با احتمال بسیار بالا یکسان با برچسب فعلی خواهد بود) با این حال، این نگرانی وجود دارد که مدل ممکن است بیش از حد به برچسب پیشین وابسته باشد و به طور کافی از توییت یا سایر اطلاعات متنی استفاده نکند. برای بررسی این موضوع، از تفسیر شپ\cite{unified_lundberg_2017}استفاده شده است تا تحلیل شود که آیا مدل از کل متن استفاده می‌کند یا به برچسب پیشین بیش از حد وابسته است. نتایج این تحلیل در بخش‌های بعدی پایان‌نامه مورد بحث قرار می‌گیرد.

\begin{table}[!ht]
\centering
\caption{ارزیابی حذف ویژگی}
\label{tbl:ablation_study}
\renewcommand{\arraystretch}{1.5} % Increase row height
\begin{tabular}{|c|c|c|c|c|} 
\hline
مدل               & صحت            & دقت            & پوشش           & F1              \\ 
\hline
بدون \lr{ROC}          & \lr{90.04}          & \lr{90.86}          & \lr{90.28}          & \lr{90.53}           \\ 
\hline
بدون \lr{RSI}          & \lr{90.03}          & \lr{90.74}          & \lr{90.38}          & \lr{90.42}           \\ 
\hline
~بدون برچسب پیشین & \textbf{\lr{43.65}} & \textbf{\lr{43.89}} & \textbf{\lr{44.07}} & \textbf{\lr{43.72}}  \\
\hline
\end{tabular}
\end{table}


\subsection{نتایج انسداد لایه}

در این بخش، نتایج انسداد لایه‌ها در مدل \lr{RoBERTa} \cite{liu2019roberta}بررسی شده است.

انسداد لایه‌ها\LTRfootnote{Freezing Layers} به معنای ثابت نگه داشتن برخی از لایه‌های مدل پیش‌آموزش‌دیده\LTRfootnote{Pre-trained} در طول فرآیند تنظیم دقیق بر روی یک وظیفه جدید است. این تکنیک به منظور جلوگیری از تغییرات زیاد در ویژگی‌های از پیش آموخته شده استفاده می‌شود و می‌تواند در تطبیق مدل‌های پیش‌آموزش‌دیده با وظایف جدید با داده‌های محدود مفید باشد. همان‌طور که در شکل مشاهده می‌شود، مدل‌هایی که تعداد کمتری از لایه‌ها را مسدود کرده‌اند، منحنی یادگیری سریعتری دارند. همانطور که در شکل \ref{fig:frozen} قابل مشاهده است به طور خاص، مدل‌هایی با ۸ لایه مسدود شده، سریع‌تر از مدل‌هایی با ۱۱ لایه مسدود شده یاد می‌گیرند. با این حال، نتایج نشان می‌دهد که پس از دومین دوره آموزشی، مدل‌ها به تقریباً همان نقطه همگرا می‌شوند. این موضوع نشان می‌دهد که اگرچه انسداد لایه‌ها می‌تواند یادگیری اولیه را تسریع کند، اما عملکرد نهایی مدل‌ها تفاوت چندانی ندارد. همچنین، باز کردن انسداد لایه‌ها می‌تواند منجر به تغییرات در لایه‌های زیرین شود که ممکن است معنای کلمات را به طور قابل توجهی تغییر دهد. بنابراین، در باز کردن انسداد لایه‌ها باید با احتیاط عمل کرد. مسدود کردن لایه‌ها به معنای ثابت نگه داشتن وزن‌های برخی لایه‌ها در طول تنظیم دقیق است. این تکنیک به حفظ دانش پیش‌آموزش مدل کمک می‌کند در حالی که به لایه‌های باقی‌مانده اجازه می‌دهد تا به وظیفه خاص منظوره مدنظر تطبیق یابند. مسدود کردن تعداد کمتری از لایه‌ها (۸) انعطاف‌پذیری بیشتری برای مدل فراهم می‌کند تا ویژگی‌های خاص وظیفه را یاد بگیرد، در حالی که مسدود کردن تعداد بیشتری از لایه‌ها (۱۱) به حفظ بیشتر دانش پیش‌آموزش کمک می‌کند، که می‌تواند در صورت محدود بودن یا نویزی بودن داده‌های آموزشی مفید باشد. آزمایش‌های ما به دنبال یافتن تعادل مناسب بین این دو جنبه برای دستیابی به بهترین عملکرد در پیش‌بینی حرکات کوتاه‌مدت بازار بر اساس توییت‌ها بود.
\cite{lee2019would}.

\begin{figure}
    \centering
    \includegraphics[width=1\linewidth]{figures/f1_score_plot_frozen.png}
    \caption{امتیاز \lr{F1} در روند یادگیری مدل آگاه به زمینه با تعداد متفاوت لایه‌های مسدود شده}
    \label{fig:frozen}
\end{figure}


\section{نتایج طبقه‌بندی}
در این بخش، نتایج طبقه‌بندی توییت‌ها به وسیله مدل‌های مختلف شامل مدل پایه، مدل نا‌آگاه به زمینه، مدل آگاه به‌ زمینه، و مدل آگاه به زمینه زمانی ارائه می‌شود. نتایج شامل جدول عملکرد طبقه‌بندی برای دو مجموعه داده است:

\begin{itemize}
    \item دسته ۱ (توییت‌های تراز‌شده برای اعتبارسنجی متقابل): این دسته شامل حدودا ۴۰ هزار توییت‌های ترکیب‌شده بوده که به طور میانگین روزانه شامل ۱۵۰ توییت می‌شود. این توییت‌ها با روش کاهش نمونه تراز شده‌اند و برای امر یادگیری و ارزیابی استفاده شده‌اند.
    \item دسته ۲ (توییت‌های دیده‌نشده توسط مدل از سال ۲۰۱۵ تا ۲۰۲۳): این دسته از توییت‌ها شامل حدودا ۲۵ هزار توییت نمونه برداری شده از اتفاقات مهم بیتکوین از سال ۲۰۱۵ تا اوایل سال ۲۰۲۳ بدون در نظر گرفتن داده‌های سال ۲۰۲۰ است.

\end{itemize}

فرایند تنظیم دقیق مدل‌های زبانی شامل چند ابرمتغیر\LTRfootnote{hyperparameter} مهم است که در کیفیت فرایند یادگیری تاثیر چشم‌‌‌گیری دارند. دو مورد از این ‌ابرمتغیرها تعداد دوره‌\LTRfootnote{epoch} و تعداد شکن\LTRfootnote{Fold} هستند که در زمینه تنظیم دقیق مدل‌های زبانی عموما بین ۲ تا ۴ دوره و ۵ تا ۱۰ شکن انتخاب می‌شود. مفهوم دوره به تعداد دفعاتی که مدل روی یک شکن از داده تنظیم می‌شود اشاره دارد و شکن به تعداد بخش‌های مجموعه‌داده که برای یادگیری و ارزیابی شکسته شده‌اند اشاره دارد. به این روش از یادگیری و ارزیابی مدل به عنوان اعتبارسنجی متقابل\LTRfootnote{Cross Validation} اشاره می‌شود که تعادلی بین کارایی محاسباتی و ارزیابی قوی مدل ایجاد می‌کند و اطمینان می‌دهد که عملکرد مدل به شدت وابسته به هیچ زیرمجموعه خاصی از داده‌ها نیست. ما از دو تکنیک مختلف اعتبارسنجی متقابل برای مدل‌های آگاه به زمینه و نا‌اگاه به زمینه استفاده کردیم. این تفاوت به دلیل وجود اطلاعات بازار در توییت‌های آگاه به زمینه اتفاق افتاده که در صورت حضور توییت‌های با روز‌های همسان منجر به نشت دیتا خواهد شد. این دو روش به این صورت تعریف شده‌اند:

\begin{itemize}
    \item اعتبارسنجی متقابل طبقه‌بندی‌شده: این روش اطمینان می‌دهد که هر شکن توزیع مشابهی از کلاس‌ها دارد، که به ویژه برای مجموعه داده‌های نامتعادل مفید است. این روش به جلوگیری از نشت داده\LTRfootnote{Data Snooping} یا تبعیض آینده‌نگر\LTRfootnote{Look-ahead Bias} می‌کند و اطمینان می‌دهد که توییت‌های یک بازه خاص در هر دو مجموعه آموزشی و ارزیابی ظاهر نمی‌شوند. این تکنیک برای مدل بدون آگاهی از زمینه استفاده شد.
    \item اعتبارسنجی متقابل گروهی\LTRfootnote{Grouped Cross Validation}: این روش اطمینان می‌دهد که هیچ توییتی از یک روز خاص در هر دو مجموعه آموزشی و ارزیابی ظاهر نمی‌شود، که برای مدل‌های آگاه از زمینه و آگاه از زمینه زمانی بسیار مهم است. با گروه‌بندی توییت‌ها بر اساس روز، ما از افشای برچسب توسط اطلاعات ارائه شده در توییت‌ها جلوگیری می‌کنیم و یکپارچگی فرآیند اعتبارسنجی را حفظ می‌کنیم.
\end{itemize}

نتایج این تحلیل‌ها نشان می‌دهند که اضافه کردن زمینه‌های مختلف، به ویژه زمینه بازار، به بهبود عملکرد مدل در پیش‌بینی حرکات کوتاه‌مدت بازار کمک کرده است. به علاوه، مدل‌های با زمینه و زمینه زمانی در برابر داده‌های خارج از نمونه توانایی خوبی در تعمیم از خود نشان داده‌اند، که قابلیت عمومی‌سازی آن‌ها را در شرایط مختلف بازار تایید می‌کند. برای تضمین ارزیابی دقیق، از اعتبارسنجی متقاطع ۵‌تایی گروه‌بندی‌شده استفاده شد که با گروه‌بندی توییت‌ها بر اساس نزدیکی زمانی، از نشت اطلاعات جلوگیری می‌کند. آموزش مدل‌ها طی ۲ دوره برای هر تقسیم‌بندی انجام شد، زیرا آزمایش‌ها نشان دادند که عملکرد مدل‌ها پس از این نقطه تثبیت می‌شود و نشان‌دهنده همگرایی موثر آن‌ها در این حد است. مدل‌های پایه، مانند \lr{FinBERT}\cite{finbert_huang_2022} و \lr{CryptoBERT}\cite{sentiment_kulakowski_2023}، که برای طبقه‌بندی احساسات از پیش آموزش داده شده بودند، عملکرد ضعیفی در پیش‌بینی روندهای کوتاه‌مدت بازار نشان دادند و به ترتیب امتیازات \lr{F1} برابر با \lr{۱۷.۲٪} و \lr{۲۷.۷٪} بر روی دسته یک کسب کردند. این نتایج محدودیت مدل‌های مبتنی بر احساسات در درک حرکات کوتاه‌مدت بازار را برجسته می‌کند. این عملکرد ضعیف به این واقعیت بازمی‌گردد که این مدل‌ها برای تفسیر برچسب‌های احساسی مبتنی بر نشانه‌های انسانی آموزش دیده‌اند، که احساسات ذهنی را منعکس می‌کند و نه رفتار واقعی بازار. در نتیجه، مدل‌ها نمی‌توانند ویژگی‌های متنی را با نتایج مشتق‌شده از بازار مرتبط کنند، زیرا به‌طور مشخص با برچسب‌های مشتق‌شده از حرکات بازار آموزش داده نشده‌اند. نتایج جدول \ref{tbl:context_aware} و شکل \ref{fig:f1_train} نشان می‌دهند که مدل‌های آگاه به زمینه و آگاه به زمینه زمانی عملکرد بهتری نسبت به مدل پایه و مدل نا‌آگاه به زمینه دارند. در ادامه به تحلیل دقیق‌تر هر مدل می‌پردازیم:


\begin{itemize}
    \item مدل \lr{FinBERT}: این مدل از \lr{FinBERT} بدون تنظیم دقیق استفاده می‌کند، عملکرد ضعیفی در هر دو دسته داده‌ها دارد. صحت، دقت، پوشش و نمره \lr{F1} این مدل در دسته ۱ به ترتیب \lr{33.37}، \lr{34.01}، \lr{33.37} و \lr{17.26} درصد است. در دسته ۲ نیز این مقادیر به ترتیب \lr{28.23}، \lr{30.80}، \lr{32.13} و \lr{25.23} درصد است.
    
    \item مدل \lr{CryptoBERT}:
    این مدل که تنها از \lr{CryptoBERT} بدون تنظیم دقیق استفاده می‌کند، عملکرد ضعیفی در هر دو دسته داده‌ها دارد. صحت، دقت، پوشش و نمره \lr{F1} این مدل در دسته ۱ به ترتیب \lr{33.35}، \lr{32.25}، \lr{33.35} و \lr{27.71} درصد است. در دسته ۲ نیز این مقادیر به ترتیب \lr{28.23}، \lr{30.80}، \lr{32.13} و \lr{25.23} درصد است. نتایج حاصل از ماتریس‌های سردرگمی نشان می‌دهد که این مدل در تشخیص کلاس‌ها (به ویژه خنثی) دچار سردرگمی زیادی است. بیشتر پیش‌بینی‌ها به صورت تصادفی میان سه کلاس تقسیم شده‌اند، که بیانگر عدم توانایی این مدل در درک سیگنال‌های بازار است.
    
    \item مدل نا‌آگاه به زمینه (\lr{CUA}):  
    این مدل که از \lr{CryptoBERT} با تنظیم دقیق برچسب‌های سه‌مانع روز بعد بدون تنظیم پرسش استفاده می‌کند، عملکرد بهتری نسبت به مدل پایه دارد. صحت، دقت، پوشش و نمره \lr{F1} این مدل در دسته ۱ به ترتیب \lr{44.13}، \lr{44.72}، \lr{44.13} و \lr{43.01} درصد است. در دسته ۲ نیز این مقادیر به ترتیب \lr{35.35}، \lr{34.46}، \lr{35.87} و \lr{33.30} درصد است. تحلیل ماتریس‌های سردرگمی نشان می‌دهد که این مدل در مقایسه با مدل پایه بهبود داشته است، اما همچنان نرخ بالایی از اشتباهات در تشخیص کلاس خنثی وجود دارد. پیش‌بینی‌های کلاس نزولی نیز در بسیاری موارد به کلاس‌های دیگر تخصیص داده شده‌اند.  
    
    \item مدل آگاه به زمینه (\lr{CA}):  
    این مدل که از \lr{CryptoBERT} با تنظیم دقیق برچسب‌های سه‌مانع روز بعد و تنظیم پرسش‌های زمینه بازار استفاده می‌کند، عملکرد بسیار بهتری نسبت به مدل‌های قبلی دارد. صحت، دقت، پوشش و نمره \lr{F1} این مدل در دسته ۱ به ترتیب \lr{89.68}، \lr{90.32}، \lr{89.34} و \lr{89.53} درصد است. در دسته ۲ نیز این مقادیر به ترتیب \lr{81.03}، \lr{80.45}، \lr{80.23} و \lr{80.32} درصد است. بررسی ماتریس‌های سردرگمی نشان می‌دهد که این مدل در جداسازی کلاس‌ها عملکرد بسیار خوبی داشته و نرخ خطاها به میزان قابل توجهی کاهش یافته است. کلاس خنثی و نزولی به‌صورت دقیق‌تر پیش‌بینی شده‌اند، که نشان‌دهنده تأثیر مثبت افزودن زمینه بازار است.  
    
    \item مدل آگاه به زمینه زمانی (\lr{TCA}):  
    این مدل که از \lr{CryptoBERT} با تنظیم دقیق برچسب‌های سه‌مانع روز بعد و تنظیم پرسش‌های زمینه بازار و زمینه زمانی استفاده می‌کند، عملکرد خوبی دارد اما کمی ضعیف‌تر از مدل آگاه به زمینه است. صحت، دقت، پوشش و نمره \lr{F1} این مدل در دسته ۱ به ترتیب \lr{86.57}، \lr{86.68}، \lr{86.58} و \lr{86.50} درصد است. در دسته ۲ نیز این مقادیر به ترتیب \lr{71.21}، \lr{71.20}، \lr{71.80} و \lr{71.08} درصد است. تحلیل ماتریس‌های سردرگمی نشان می‌دهد که این مدل نیز قادر به پیش‌بینی دقیق کلاس‌ها است، اما اندکی خطاهای بیشتری نسبت به مدل آگاه به زمینه (\lr{CA}) دارد. به‌ویژه در کلاس خنثی تعداد کمی از نمونه‌ها به سایر کلاس‌ها تخصیص داده شده‌اند. این کاهش عملکرد ممکن است ناشی از پیچیدگی اضافی مدل در مدیریت زمینه زمانی باشد.  

\end{itemize}


\begin{table}
\centering
\caption{ارزیابی طبقه‌بندی توییت‌ها}
\label{tbl:context_aware}
\begin{tabular}{c|cc|cc|cc|cl} 
\hline
\multirow{2}{*}{مدل پیشبینی} & \multicolumn{2}{c|}{صحت} & \multicolumn{2}{c|}{دقت} & \multicolumn{2}{c|}{پوشش} & \multicolumn{2}{c}{\lr{F1}}               \\ 
                             & دسته 1 & دسته 2          & دسته 1 & دسته 2          & دسته 1 & دسته 2           & دسته 1                     & دسته 2  \\ 
\hline
\lr{FinBERT}                         
& \lr{33.37}          & \lr{32.82}          & \lr{34.01}          & \lr{26.4}          & \lr{33.37}          & \lr{33.24}          & \lr{17.26}          & \lr{16.7}\\
\lr{CryptoBERT}                         & \lr{33.35}  & \lr{28.23}           & \lr{32.25}  & \lr{30.80}           & \lr{33.35}  & \lr{32.13}            & \lr{27.71} & \lr{25.23}   \\
مدل نا‌آگاه به زمینه (\lr{CUA})        & \lr{44.13}  & \lr{35.35}           & \lr{44.72}  & \lr{34.46}           & \lr{44.13}  & \lr{35.87}            & \lr{43.01}                      & \lr{33.30}   \\
مدل آگاه به زمینه (\lr{CA})           & \lr{89.68}  & \lr{81.03}           & \lr{90.32}  & \lr{80.45}           & \lr{89.34}  & \lr{80.23}            & \multicolumn{1}{c:}{\lr{89.53}} & \lr{80.32}   \\
مدل آگاه به زمینه زمانی (\lr{TCA})      & \lr{86.57}  & \lr{71.21}           & \lr{86.68}  & \lr{71.20}           & \lr{86.58}  & \lr{71.80}            & \lr{86.50}                      & \lr{71.08}   \\
\hline
\end{tabular}
\end{table}


\begin{figure}
    \centering
    \includegraphics[width=1\linewidth]{figures/classification_confusion_matrix.png}
    \caption{ماتریس در‌هم‌ریختگی مدل‌های مختلف برای طبقه‌بندی توییت‌ها}
    \label{fig:f1_train}
\end{figure}

مدل‌های \lr{CA} و \lr{TCA} اگرچه امتیازات \lr{F1} بالاتری کسب کردند، اما نوسانات بیشتری نیز نشان دادند. به‌عنوان مثال، امتیازات \lr{F1} مدل \lr{TCA} در محدوده \lr{81.39} تا \lr{92.20} درصد (\lr{std = 0.0399}) و مدل \lr{CA} در محدوده \lr{85.28} تا \lr{96.48} درصد (\lr{std = 0.0413}) بود. برای کاهش این نوسانات، از استراتژی رأی‌گیری اکثریت\LTRfootnote{Majority Voting} استفاده شد که نوعی یادگیری جمعی\LTRfootnote{Ensemble Learning} به حساب می‌آید. در این رویکرد، پیش‌بینی‌های مدل‌های آموزش‌دیده روی تقسیم‌بندی‌های مختلف داده‌ها ترکیب شدند تا تصمیم نهایی درباره دسته‌بندی در مجموعه داده‌های دیده‌نشده اتخاذ شود. این روش که مبتنی بر تجمیع پیش‌بینی‌ها در پنج تقسیم‌بندی است، با کاهش تأثیر ناهنجاری‌های خاص هر تقسیم‌بندی داده، ثبات پیش‌بینی‌ها را افزایش داد و ارزیابی قابل‌اعتمادتری از عملکرد مدل ارائه کرد. علاوه بر این، \lr{paired t-test} برای مقایسه مدل‌های \lr{CA} و \lr{TCA} با \lr{CryptoBERT} در پنج تقسیم‌بندی، تفاوت‌های آماری معنادار در عملکرد را تأیید کردند (\lr{$p < 0.001$} برای هر دو مقایسه). این نتایج برتری مدل‌های پیشنهادی نسبت به \lr{CryptoBERT} را تقویت می‌کند و نشان می‌دهد که بهبود عملکرد آن‌ها به‌دلیل تغییرات تصادفی داده‌ها نیست.

\begin{figure*}[ht!]
    \centering
    \includegraphics[width=\textwidth]{figures/shap.png}
    \caption{نمودار مقادیر شپ برای توییت. این نمودار تأثیر هر کلمه بر پیش‌بینی مدل را نشان می‌دهد و تأثیر قابل توجه متن توییت بر پیش‌بینی نزولی را برجسته می‌کند.}
    \label{fig:shap}
\end{figure*}


نتایج آزمایش‌های ما نکات ارزشمندی را در مورد قدرت پیش‌بینی مدل‌های پیشنهادی و تأثیرات آن‌ها بر پیش‌بینی بازارهای مالی فراهم می‌آورد. یکی از نگرانی‌های اصلی احتمال بیش‌برازش مدل بر روی پرسش‌ها بود، به طوری که مدل ممکن است تنها به پرسش‌های در برچسب‌گذاری توییت‌ها تکیه کند. این می‌تواند منجر به آن شود که مدل تنها بر اساس برچسب قبلی عمل کند به دلیل همبستگی بالای آن با پیش‌بینی بعدی. برای رفع این نگرانی، نتایج با استفاده از مقادیر شپ تفسیر شد. نمودار مقادیر شپ (شکل \ref{fig:shap}) برای توییت انتخاب‌شده تصویری واضح از نحوه پردازش اطلاعات توسط مدل ارائه می‌دهد. با وجود برچسب قبلی صعودی، مدل پیش‌بینی نزولی می‌کند. این عمدتاً به دلیل تأثیر قابل توجه کلمه "downward" از محتوای توییت است. مقادیر شپ نشان می‌دهند که مدل به‌طور مؤثر محتوای توییت را تجزیه و تحلیل می‌کند و تنها به پرسش‌ها تکیه نمی‌کند. کلمه "downward" دارای مقدار شپ مثبت و قابل توجهی است که به شدت بر پیش‌بینی نهایی مدل به سمت نزولی تأثیر می‌گذارد. این نشان می‌دهد که مدل محتوای واقعی توییت را در نظر می‌گیرد و پیش‌بینی‌ها بر اساس تحلیل معنادار و نه اطلاعات سطحی پرسش‌ها است.

اهمیت این یافته بر آن است که استحکام مدل ما در تفسیر توییت‌ها برای پیش‌بینی بازار با استفاده‌ از اطلاعات زمینه‌ای استفاده شده را تأیید می‌کند. نشان می‌دهد که مدل می‌تواند اطلاعات مرتبط را از محتوای توییت شناسایی کند، حتی زمانی که پرسش‌ها احساسات متفاوتی را پیشنهاد می‌دهند. این قابلیت قدرت پیش‌بینی مدل را افزایش می‌دهد و آن را به ابزاری ارزشمند برای تحلیل بازارهای مالی تبدیل می‌کند. شواهد تجربی از آزمایش‌های ما نشان می‌دهد که روش سه‌مانع، همراه با مهندسی پرسش مؤثر، تحلیل دقیق‌تر و معنادارتری از محتوای رسانه‌های اجتماعی مالی ارائه می‌دهد. این نکته به‌طور کلی نقاط قوت رویکرد مدل‌سازی ما را برجسته می‌کند و راهنمایی‌های ارزشمندی برای بهبودها و کاربردهای آینده در پیش‌بینی بازارهای مالی ارائه می‌دهد. با تمرکز بر تأثیر واقعی بازار توییت‌ها، می‌توانیم تحلیل دقیق‌تر و معنادارتری ارائه دهیم و قدرت پیش‌بینی کلی مدل‌های خود را افزایش دهیم.

نتایج این تحلیل‌ها نشان می‌دهند که اضافه کردن اطلاعات زمینه‌‌ای مختلف، به ویژه زمینه بازار، به بهبود عملکرد مدل در پیش‌بینی حرکات کوتاه‌مدت بازار کمک کرده است. به علاوه، مدل‌های با زمینه و زمینه زمانی در برابر داده‌های خارج از نمونه توانایی خوبی در تعمیم از خود نشان داده‌اند، که قابلیت عمومی‌سازی آن‌ها را در شرایط مختلف بازار تایید می‌کند.


\newpage
\section{تولید سیگنال}
فرآیند تولید سیگنال، طبقه‌بندی‌های توییت‌ها را به سیگنال‌های معاملاتی قابل‌اجرا تبدیل می‌کند و توصیه‌ای نهایی و روشن در مورد حرکت کوتاه‌مدت بازار ارائه می‌دهد. برخلاف وظیفه‌ی طبقه‌بندی که بر روی توییت‌های منفرد تمرکز دارد، این مرحله پیش‌بینی‌های تمام توییت‌های یک روز معین را برای تعیین سیگنال غالب تجمیع می‌کند. روش‌های سنتی تولید سیگنال اغلب به معماری‌های چندمرحله‌ای متکی هستند، جایی که پیش‌بینی‌های احساسات با ویژگی‌های مشتق‌شده از بازار در یک مدل ترکیبی جداگانه ترکیب می‌شوند تا پیش‌بینی نهایی از حرکت بازار را ارائه دهند. اگرچه این روش‌ها در برخی موارد مؤثر هستند، اما چنین معماری‌هایی نمی‌توانند تعاملات ظریف بین ویژگی‌های متنی و اطلاعات زمینه‌ای بازار را به‌ خوبی استخراج کنند، که این امر منجر به از دست رفتن اطلاعات می‌شود و به داده‌های بیشتر و معماری پیچیده‌تری نیاز دارند. نتایج آزمایش تولید سیگنال در جدول \ref{tbl:next_day} و شکل \ref{fig:true_p} خلاصه شده است. هدف اصلی این آزمایش ارزیابی کارایی مدل در تولید سیگنال‌های معاملاتی قابل استفاده در بازار‌های مالی بر اساس پیشبینی‌های انجام شده بر توییت‌ها است. عملکرد مدل حافظه بلند کوتاه‌مدت با دو روش دیگر، روش رأی‌گیری اکثریت و روش رأی‌گیری میانگین برای هر روز مقایسه شده که نتایج به دست آمده به طور کامل در جدول \ref{tbl:next_day}.

در مقابل، مدل زبانی \lr{CA} ویژگی‌های بازاری مانند روند پیشین، نرخ تغییر و شاخص قدرت نسبی را مستقیماً در قالب توصیف‌های متنی به پرسش توییت‌ها اضافه می‌کند. این رویکرد امکان یک فرآیند پیش‌بینی انتها به انتها را فراهم می‌کند که به‌صورت ذاتی این تعاملات را مدل‌سازی می‌کند. این مطالعه اثربخشی این رویکرد نوآورانه را در تولید سیگنال‌های قابل‌اطمینان ارزیابی کرده و عملکرد آن را با مدل‌های سنتی مبتنی بر قیمت و مدل‌های ترکیبی مقایسه می‌کند. برای ارزیابی این روش، دو مدل مبتنی بر قیمت پیاده‌سازی شدند. مدل اول یک \lr{LSTM} است که از \cite{prebit_zou_2023} الهام گرفته و مدل دوم یک \lr{Autoencoder} است که از \cite{zha2022time} الهام گرفته شده است. هر دو مدل بر روی داده‌های سال‌های ۲۰۱۵ تا ۲۰۱۹ آموزش دیده‌اند و از ویژگی‌هایی مانند داده‌های شمعی\LTRfootnote{OHLCV} بیتکوین، قیمت اتریوم، قیمت طلا، \lr{RSI}، \lr{ROC}، \lr{MACD} و میانگین‌های متحرک (\lr{MA7}، \lr{MA21}، \lr{EMA12} و \lr{EMA26}) استفاده کردند. علاوه بر این، یک مدل ترکیبی نیز مورد بررسی قرار گرفت که پیش‌بینی‌های مدل \lr{CA}، پیش‌بینی‌های احساسات و پیش‌بینی‌های مبتنی بر قیمت را برای تولید سیگنال‌های روزانه ترکیب می‌کند. پیش‌بینی‌های احساسات به‌صورت نسبت توییت‌های صعودی به نزولی در روز ارائه شدند، درحالی‌که پیش‌بینی‌های مبتنی بر قیمت از مدل‌های \lr{LSTM} و \lr{Autoencoder} به مدل ترکیبی داده شدند. در نهایت، مدل \lr{CA} پیش‌بینی‌های طبقه‌بندی‌شده‌ی توییت‌ها را با استفاده از رأی‌گیری اکثریت در یک تنظیم ارزیابی متقابل روزانه تجمیع کرد. برخلاف مدل‌های مبتنی بر قیمت که به چهار سال داده‌ی آموزشی نیاز داشتند، مدل \lr{CA} با استفاده از تنها یک سال توییت (۲۰۲۰) به نتایج قابل‌مقایسه‌ای دست یافت که کارآمدی آن را نشان می‌دهد.

\begin{table}
\centering
\renewcommand{\arraystretch}{1.5} % Increase row height
% \setlength\extrarowheight{5pt}
\caption{ارزیابی سیگنال‌های روزانه}
\label{tbl:next_day}
\begin{tabular}{cccccc}
\multicolumn{2}{c}{مدل پیشبینی}           & صحت            & دقت            & پوشش           & F1              \\ 
\hline
\multirow{3}{*}{تمام دسته‌ها}   & 
                                \lr{LSTM (Price)}
                                & \lr{52.87}
                                & \lr{46.76}
                                & \lr{52.87}
                                & \lr{49.20}          \\
                                & \lr{AutoEncoder (Price)}    
                                & \lr{54.25}
                                & \lr{51.31}
                                & \lr{54.63}
                                & \lr{50.01}          \\

                                & \lr{Fusion}    & \textbf{\lr{59.01}} & \lr{58.35} & \lr{55.14}          & \textbf{\lr{55.94}} \\
                                & \lr{اکثریت}  &
                                \lr{53.52}          & \textbf{\lr{61.15}}          & \textbf{\lr{57.46}} & \lr{54.21}          \\
                                & \lr{میانگین} &
                                \lr{57.91}          & \lr{56.50}          & \lr{57.13}          & \lr{56.10}          \\ 
\hline
\multirow{3}{*}{صعودی به نزولی} & \lr{Fusion}    &
\lr{78.38}          & \lr{80.72}          & \textbf{\lr{90.38}} & \lr{85.28} \\
                                & \lr{اکثریت}  & \textbf{\lr{84.57}} & \textbf{\lr{84.30}} & \lr{88.02} & \textbf{\lr{86.12}}          \\
                                & \lr{میانگین} & \lr{76.95}          & \lr{83.60}          & \lr{76.33}          & \lr{79.80}          \\
\hline
\multirow{3}{*}{نزولی به صعودی} & \lr{Fusion}    & \lr{78.38}          & \lr{70.31}          & \lr{51.33}          & \lr{59.34}          \\
                                & \lr{اکثریت}  & \textbf{\lr{84.92}} & \textbf{\lr{84.92}} & \textbf{\lr{80.45}} & \textbf{\lr{82.62}} \\
                                & \lr{میانگین} & \lr{76.95}          & \lr{69.00}          & \lr{77.86}          & \lr{73.16}                  
\end{tabular}
\end{table}

برای ارزیابی عملکرد این مدل‌ها، معیارهای صحت، دقت، یادآوری و امتیاز \lr{F1} در دو حالت یک در برابر همه \LTRfootnote{One-Vs-Rest} و یک در برابر یک \LTRfootnote{One-Vs-One} استفاده شدند. در حالت یک در برابر همه، توانایی مدل‌ها برای تمایز بین سه کلاس صعودی، نزولی و خنثی مورد بررسی قرار گرفت. در حالت یک در برابر یک، مقایسه‌های جفتی، مانند صعودی در مقابل نزولی، ارزیابی شد که اهمیت بیشتری برای تصمیم‌گیری‌های معاملاتی دارند. این ارزیابی جامع، توانایی مدل‌ها را به‌خوبی نشان داد و حساسیت روش یک در برابر یک به سناریوهای مهم تصمیم‌گیری را برجسته کرد. نتایج، که در جدول \ref{tbl:next_day} خلاصه شده‌اند، نشان می‌دهند که مدل \lr{CA} عملکرد رقابتی و قابل‌توجهی ارائه می‌دهد. این مدل با استفاده از رأی‌گیری اکثریت به امتیاز کلی \lr{F1} معادل \lr{54.21\%} دست یافت و مدل‌های (\lr{LSTM} = \lr{49.20\%}) و (\lr{Autoencoder} = \lr{50.01\%}) را پشت سر گذاشت. همچنین در کلاس‌های جداگانه، عملکرد این مدل برتر بود. به‌عنوان مثال، امتیاز \lr{F1} برای برچسب صعودی برابر با \lr{86.12\%} بود، درحالی‌که مدل ترکیبی امتیاز \lr{85.28\%} و مدل \lr{LSTM} امتیاز \lr{79.34\%} داشتند. برای برچسب نزولی نیز مدل \lr{CA} با امتیاز \lr{82.62\%} عملکرد بهتری نسبت به مدل ترکیبی \lr{59.34\%} و مدل (\lr{64.21\%} = \lr{LSTM}) داشت. یکی از مشاهدات مهم از ماتریس اغتشاش روش رأی‌گیری اکثریت، رفتار محافظه‌کارانه‌ی آن در خطاسازی است. تنها \lr{5.1\%} از برچسب‌های صعودی به‌عنوان نزولی و \lr{12.5\%} از برچسب‌های نزولی به‌عنوان صعودی خطاسازی شدند. اکثر خطاها شامل پیش‌بینی برچسب‌های صعودی یا نزولی به‌عنوان خنثی بودند که این موضوع خطرات معاملاتی را با جلوگیری از سیگنال‌های بیش‌ازحد تهاجمی در سناریوهای نامطمئن کاهش می‌دهد.

\begin{figure}
    \centering
    \includegraphics[width=1\linewidth]{figures/precision_recall_scatter.png}
    \caption{نسبت مثبت واقعی به نسبت مثبت کاذب برای دو روش اکثریت و میانگین در پیشبینی برچسب روز بعد}
    \label{fig:true_p}
\end{figure}

 شکل \ref{fig:true_p} مقایسه‌ای از معیارهای دقت-بازخوانی\LTRfootnote{Precision-Recall}برای روش‌های مختلف تجمیع توییت‌ها به منظور پیش‌بینی احساسات بازار را نشان می‌دهد. این احساسات در دو دسته  و صعودی طبقه‌بندی شده‌اند. هر نقطه داده نشان‌دهنده یک روش خاص است: روش‌های تجمیع اکثریت و میانگین به صورت جداگانه استفاده شده‌اند، در حالی که روش ترکیبی\LTRfootnote{Fusion} هر دو را با هم ترکیب می‌کند. امتیازات \lr{F1} برای هر کلاس احساسات در کنار نقاط مربوطه نشان داده شده است و عملکرد روش‌ها را به تصویر می‌کشد. برای کلاس صعودی، روش تجمیع اکثریت به بالاترین امتیاز \lr{F1} برابر با \lr{86.01\%} دست یافته است و پس از آن روش ترکیبی با امتیاز \lr{85.3} قرار دارد. برای کلاس نزولی، روش اکثریت نیز عملکرد قوی‌ای با امتیاز \lr{F1} برابر با \lr{82.6\%} نشان می‌دهد، در حالی که روش ترکیبی کمترین امتیاز \lr{F1} برابر با \lr{59.3\%} را دارد. روش میانگین نتایج میانی را با امتیاز \lr{79.8\%} برای صعودی و \lr{73.2\%} برای نزولی ارائه می‌دهد. این نمودار پویایی‌های دقت-بازخوانی را نشان می‌دهد و بر نقاط قوت و ضعف هر استراتژی تجمیع در مدیریت پیش‌بینی‌های احساسات صعودی و نزولی تأکید می‌کند. در نهایت، نتایج ارائه‌شده در شکل نشان می‌دهند که مدل \lr{CA} توانایی حفظ دقت و یادآوری بالا را دارد و درعین‌حال از خطاهای پرریسک اجتناب می‌کند. این یافته‌ها پتانسیل مدل \lr{CA} را به‌عنوان یک رویکرد انتها به انتها برای تولید سیگنال‌های معاملاتی قابل‌اطمینان برجسته می‌کنند. با ادغام بافت بازاری در پرسش‌های توییت و دور زدن نیاز به معماری‌های ترکیبی جداگانه، مدل \lr{CA} نشان می‌دهد که چگونه مدل‌های زبانی می‌توانند تعاملات پیچیده بین داده‌های متنی و بازاری را به‌طور مؤثری ثبت کنند. این نوآوری اتکا به معماری‌های چندمرحله‌ای را کاهش داده و جایگزینی ساده و مؤثر برای پیش‌بینی مالی ارائه می‌دهد.


\section{نتایج پیش‌آزمون}
در این بخش، پس از تولید سیگنال‌ها، عملکرد روش‌های پیشنهادی را از نظر توانایی آن‌ها در کسب سود از بازار با استفاده از چند استراتژی ساده معاملاتی ارزیابی می‌کنیم. یکی از چالش‌های مهم در ارزیابی استراتژی‌های معاملاتی، وابستگی آن‌ها به رژیم‌های مختلف بازار است که می‌تواند بر عملکرد تأثیر بگذارد. برای ارزیابی جامع، استراتژی‌های ما در سه رژیم متفاوت بازار آزمایش شدند: رژیم صعودی از سال ۲۰۱۷ تا ۲۰۱۸ که با رشد سریع بازار به میزان تقریباً \lr{1200\%} مشخص می‌شود؛ رژیم نزولی از سال ۲۰۱۸ تا ۲۰۱۹ که با کاهش شدید \lr{250\%} همراه بود؛ و رژیم خنثی از اواسط سال ۲۰۱۹ تا اواسط سال ۲۰۲۰ که ثبات نسبی با افزایش حدود \lr{40\%} را نشان داد. روش‌های تولید سیگنال پیشنهادی در کنار سه استراتژی معاملاتی اعمال شدند. برای مقایسه، عملکرد این روش‌ها با معیارهای خرید و نگهداری و فروش و نگهداری مقایسه شد که به عنوان پایه‌ای برای استراتژی‌های منفعل که روندهای بازار را دنبال می‌کنند، عمل می‌کنند. با آزمایش این استراتژی‌ها در شرایط مختلف بازار، هدف ما درک سازگاری و استحکام روش‌های پیشنهادی در تولید سیگنال‌های معاملاتی قابل اجرا بود.

نتایج ارائه‌شده در جدول \ref{tbl:backtest_results} الگوهای قابل‌توجهی را برجسته می‌کنند. در دوره صعودی، استراتژی معیار خرید و نگهداری به نسبت شارپ \lr{3.14} و بازدهی روزانه \lr{3.38\%}  دست یافت که روند صعودی فوق‌العاده بیت‌کوین را منعکس می‌کند. در حالی که عملکرد معیار به‌طور استثنایی خوب بود، چهار مورد از شش استراتژی پیشنهادی از این نسبت شارپ فراتر رفتند و استراتژی اکثریت سه‌مانع به مقدار \lr{10.5}رسید. در دو رژیم دیگر بازار، همه استراتژی‌های پیشنهادی از هر دو معیار خرید و نگهداری و فروش و نگهداری\LTRfootnote{Sell and Hold} عملکرد بهتری داشتند و سازگاری قوی خود را نشان دادند. استراتژی‌های ورود و خروج وابستگی زیادی به رژیم بازار داشتند، همان‌طور که انتظار می‌رود با توجه به رویکرد معاملاتی یک‌طرفه آن‌ها. به عنوان مثال، استراتژی ورود و خروج صعودی در رژیم صعودی عملکرد فوق‌العاده‌ای داشت و بازده روزانه بالایی ارائه داد، اما در شرایط نزولی عملکرد کمتری داشت. این استراتژی‌ها مدت‌زمان معاملات طولانی‌تری نسبت به استراتژی‌های سه‌مانع داشتند که منجر به افزایش ریسک شد، و معاملات کمتری اجرا کردند، اما نسبت شارپ قابل قبولی داشتند—به‌جز زمانی که برخلاف روند بازار عمل می‌کردند.


\begin{table}
\centering
\caption{ارزیابی استراتژی‌ها در پیش‌آزمون}
\label{tbl:backtest_results}
\begin{tabular}{lcccccccc}
                                 & \multicolumn{2}{c}{استراتژی}                 & بازدهی        & شارپ~         & سورتینو       & معاملات      & افت سرمایه     & نسبت سود        \\ 
\hline
\multirow{7}{*}{\rotcell{صعودی}} & \multicolumn{2}{c}{خرید و نگه‌داری}          & \lr{3.38}     & \lr{3.14}     & \lr{4.55}     & \lr{1}       & \lr{35.36}     & -               \\ 
\cline{2-9}
                                 & \multirow{3}{*}{اکثریت}  & ورود و خروج صعودی & \lr{10.88}    & \lr{4.86}     & \lr{7.18}     & \lr{11}      & \lr{34.46}     & \lr{807.90}  \\
                                 &                          & ورود و خروج نزولی & \lr{0.40}     & \lr{2.44}     & \lr{2.54}     & \lr{11}      & \lr{11.11}     & \lr{19.44}           \\
                                 &                          & سه‌مانع           & \lr{7.86}     & \lr{5.10}     & \lr{8.47}     & \lr{66}      & \lr{20.92}     & \lr{12.04}\\ 
\cline{2-9}
                                 & \multirow{3}{*}{میانگین} & ورود و خروج صعودی & \lr{6.45}     & \lr{4.61}     & \lr{6.02}     & \lr{17}      & \lr{21.56}     & \lr{290.19}     \\
                                 &                          & ورود و خروج نزولی & \lr{0.01}     & \lr{0.39}     & \lr{0.40}     & \lr{16}      & \lr{61.01}     & \lr{1.04}       \\
                                 &                          & سه‌مانع           & \lr{4.27}     & \lr{3.72}     & \lr{6.16}     & \lr{79}      & \lr{47.04}     & \lr{2.51}       \\ 
\hline
\multirow{7}{*}{\rotcell{خنثی}} & \multicolumn{2}{c}{خرید و نگه‌داری}          & \lr{3.38}     & \lr{3.14}     & \lr{4.55}     & \lr{1}       & \lr{35.36}     & -               \\ 
\cline{2-9}
                                 & \multirow{3}{*}{اکثریت}  & ورود و خروج صعودی & \lr{10.88}    & \lr{4.86}     & \lr{7.18}     & \lr{11}      & \lr{34.46}     & \lr{807.90}     \\
                                 &                          & ورود و خروج نزولی & \lr{0.40}     & \lr{2.44}     & \lr{2.54}     & \lr{11}      & \lr{11.11}     & \lr{19.44}      \\
                                 &                          & سه‌مانع           & \lr{7.86}     & \lr{5.10}     & \lr{8.47}     & \lr{66}      & \lr{20.92}     & \lr{12.04}      \\ 
\cline{2-9}
                                 & \multirow{3}{*}{میانگین} & ورود و خروج صعودی & \lr{6.45}     & \lr{4.61}     & \lr{6.02}     & \lr{17}      & \lr{21.56}     & \lr{290.19}     \\
                                 &                          & ورود و خروج نزولی & \lr{0.01}     & \lr{0.39}     & \lr{0.40}     & \lr{16}      & \lr{61.01}     & \lr{1.04}       \\
                                 &                          & سه‌مانع           & \lr{4.27}     & \lr{3.72}     & \lr{6.16}     & \lr{79}      & \lr{47.04}     & \lr{2.51}       \\ 
\hline
\multirow{7}{*}{\rotcell{نزولی}} & \multicolumn{2}{c}{خرید و نگه‌داری}          & \lr{3.38}     & \lr{3.14}     & \lr{4.55}     & \lr{1}       & \lr{35.36}     & -               \\ 
\cline{2-9}
                                 & \multirow{3}{*}{اکثریت}  & ورود و خروج صعودی & \lr{10.88}    & \lr{4.86}     & \lr{7.18}     & \lr{11}      & \lr{34.46}     & \lr{807.90}     \\
                                 &                          & ورود و خروج نزولی & \lr{0.40}     & \lr{2.44}     & \lr{2.54}     & \lr{11}      & \lr{11.11}     & \lr{19.44}      \\
                                 &                          & سه‌مانع           & \lr{7.86}     & \lr{5.10}     & \lr{8.47}     & \lr{66}      & \lr{20.92}     & \lr{12.04}      \\ 
\cline{2-9}
                                 & \multirow{3}{*}{میانگین} & ورود و خروج صعودی & \lr{6.45}     & \lr{4.61}     & \lr{6.02}     & \lr{17}      & \lr{21.56}     & \lr{290.19}     \\
                                 &                          & ورود و خروج نزولی & \lr{0.01}     & \lr{0.39}     & \lr{0.40}     & \lr{16}      & \lr{61.01}     & \lr{1.04}       \\
                                 &                          & سه‌مانع           & \lr{4.27}     & \lr{3.72}     & \lr{6.16}     & \lr{79}      & \lr{47.04}     & \lr{2.51}       \\
\hline
\end{tabular}
\end{table}

در میان روش‌های پیشنهادی، استراتژی‌های سه‌مانع به‌عنوان برترین عملکردها ظاهر شدند، با مدت‌زمان معاملات کوتاه‌تر (معمولاً ۲ تا ۵ روز)، کاهش حداکثر افت سرمایه، و کاهش ریسک. این استراتژی‌ها معاملات بیشتری نسبت به استراتژی‌های ورود و خروج اجرا کردند و به‌طور مداوم نسبت‌های شارپ و سورتینو قوی ارائه دادند، و عملکرد پایداری را در رژیم‌های مختلف، از جمله سودآوری در بازارهای بی‌حرکت خنثی، نشان دادند. در مورد روش‌های تجمیع، در حالی که رویکردهای اکثریت و میانگین نتایج کلی مشابهی داشتند، روش اکثریت در بازارهای رونددار (صعودی یا نزولی) برتری داشت، در حالی که روش میانگین عملکرد بهتری در بازارهای خنثی نشان داد. 

در نهایت، آزمایش استراتژی‌های اکثریت سه‌مانع و میانگین سه‌مانع در ژانویه ۲۰۲۱—اولین ماه پس از دوره آموزش و یک فاز صعودی—نسبت‌های شارپ \lr{4.43} و \lr{3.49} به ترتیب به‌دست آورد که با عملکرد مورد انتظار آن‌ها در چنین شرایطی همخوانی داشت. این نتایج کاربردپذیری و استحکام پیش‌بینی‌های مبتنی بر مدل زبان ما را در محیط‌های پیش‌آزمون برجسته می‌کند. فرآیند تنظیم دقیق، که به‌طور خاص با برچسب‌های مشتق‌شده از بازار سه‌مانع طراحی شده بود، به‌طور طبیعی نقاط قوت مدل را با استراتژی سه‌مانع هماهنگ کرد. این نشان می‌دهد که انتخاب روش برچسب‌گذاری می‌تواند به‌عنوان یک مزیت استراتژیک عمل کند و امکان استفاده از استراتژی‌هایی که به‌طور مستقیم با فرآیند برچسب‌گذاری همسو هستند را فراهم آورد تا به نتایج برتری دست یابند. استراتژی سه‌مانع به‌عنوان یک رویکرد قابل‌اعتماد برای استفاده از پیش‌بینی‌های مدل برجسته می‌شود و یک پروفایل ریسک متعادل و سازگاری با شرایط متنوع بازار ارائه می‌دهد.


استراتژی برچسب‌گذاری سه‌مانع در بازده‌های منظم‌تری از دیگران پیشی گرفت و بالاترین نسبت شارپ را کسب کرد. این استراتژی مدیریت ریسک مؤثری را با مقدار متعادلی از حداکثر افت سرمایه و تعداد بالای معاملات بسته‌شده نشان داد و ثبات و قابل اعتماد بودن آن را در پیش‌بینی حرکت‌های بازار به نمایش گذاشت. این یافته‌ها اهمیت در نظر گرفتن بازده و ریسک هنگام ارزیابی استراتژی‌های تجاری را برجسته می‌کند. عملکرد برتر استراتژی سه‌مانع از نظر بازده‌های با ثبات‌تر، آن را به یک رویکرد امیدوارکننده برای پیش‌بینی و معامله در بازارهای مالی تبدیل می‌کند. نتایج همچنین پتانسیل ادغام مدل‌های زبان پیشرفته و اطلاعات زمینه‌ای برای افزایش دقت پیش‌بینی در بازارهای مالی را نشان می‌دهد.

\subsection{تحلیل مقایسه‌ای بین پرتفوی روش معیار و روش پیشنهادی}

\begin{figure}
    \centering
    \includegraphics[width=1\linewidth]{figures/backtest.png}
    \caption{استراتژی سه‌مانع و معاملات انجام شده در بازه  ۲۰۱۹-۲۰۲۰}
    \label{fig:enter-label}
\end{figure}

نتایج پیش‌آزمون به وضوح رفتار متفاوت پرتفوی\LTRfootnote{Portfolio} روش پیشنهادی و روش معیار را در دوره زمانی سال‌های ۲۰۱۹ تا ۲۰۲۰ نشان می‌دهد. در ابتدای سال، بازار روند صعودی قدرتمندی را تجربه می‌کند که این موضوع در عملکرد پرتفوی روش معیار (خط مشکی در نمودار بازده تجمعی) منعکس شده است. از آنجایی که روش معیار رویکرد خرید و نگه‌داری را دنبال می‌کند، ارزش آن به همراه افزایش قیمت‌های بازار افزایش می‌یابد. در ابتدا، روش معیار به دلیل اتکای کامل به روند کلی بازار، عملکردی کمی بهتر از روش پیشنهادی (خط بنفش) دارد، زیرا روش پیشنهادی با احتیاط بیشتری به اجرای معاملات می‌پردازد و به جای اتکا به حرکت کلی صعودی بازار، بر فرصت‌های کسب سود از نوسانات قیمت تمرکز دارد.

با این حال، تفاوت میان دو پرتفوی هنگامی که شرایط بازار در میانه سال تغییر می‌کند، آشکار می‌شود. ارزش پرتفوی روش معیار به طور قابل توجهی در واکنش به آشفتگی‌های بازار کاهش می‌یابد و آسیب‌پذیری آن در برابر نزول بازار را نشان می‌دهد. در مقابل، روش پیشنهادی توانایی خود را در تطبیق با شرایط مختلف بازار به نمایش می‌گذارد. این روش با استفاده از سیگنال‌های خرید و فروش به موقع، از حرکت‌های کوتاه‌مدت قیمت بهره می‌برد؛ شواهد آن را می‌توان در تجمع معاملات سودآور (دایره‌های سبز) در دوره‌های پرتلاطم بازار مشاهده کرد. این توانایی در تنظیم پوزیشن‌ها به صورت پویا باعث می‌شود روش پیشنهادی حتی در زمانی که روش معیار با کاهش قیمت‌ها دست و پنجه نرم می‌کند، به روند صعودی بازده تجمعی خود ادامه دهد. بررسی دقیق‌تر نمودار دوم، که سود و زیان معاملات را نمایش می‌دهد، قدرت روش پیشنهادی را بیشتر نشان می‌دهد. این روش در هر دو شرایط صعودی و نزولی بازار به طور مداوم معاملات خود را با سود می‌بندد. این موضوع به خصوص در دوره‌های پرتلاطم بازار مشهود است؛ جایی که نقاط ورود و خروج دقیق، زیان‌ها را کاهش داده و فرصت‌های کسب سود را شناسایی می‌کند. در مقابل، روش معیار که به روند کلی بازار وابسته است، فاقد چنین انعطافی است و از نزول‌های بازار آسیب می‌بیند.

تا پایان دوره آزمایش، روش پیشنهادی به طور قطعی عملکرد بهتری نسبت به روش معیار نشان می‌دهد. در حالی که عملکرد روش معیار در مواجهه با شرایط نامساعد بازار دچار رکود یا حتی کاهش می‌شود، روش پیشنهادی همچنان به رشد خود ادامه می‌دهد و توانایی آن در کسب سود از رژیم‌های مختلف بازار را نشان می‌دهد. این تطبیق‌پذیری یک مزیت حیاتی است؛ چرا که به استراتژی اجازه می‌دهد نه تنها در بازارهای دارای روند بلکه در دوره‌های نوسانات و بازگشت‌ها نیز عملکرد قوی خود را حفظ کند. به طور خلاصه، در حالی که پرتفوی روش معیار در ابتدا از روند صعودی بازار بهره می‌برد، وابستگی سختگیرانه آن به استراتژی خرید و نگه‌داری، عملکرد آن را در دوره‌های پرتلاطم محدود می‌کند. در مقابل، روش پیشنهادی با انعطاف‌پذیری و مقاومت بالاتر خود، به سرعت از روش معیار پیشی می‌گیرد و با موفقیت شرایط متغیر بازار را مدیریت می‌کند. این تفاوت، قدرت روش پیشنهادی را نشان می‌دهد، به ویژه در توانایی آن در کسب سود از آشفتگی‌ها و رژیم‌های مختلف بازار که آن را به انتخابی مطمئن‌تر برای رشد پایدار پرتفوی تبدیل می‌کند.


\section*{نتیجه‌گیری}

یافته‌های ارائه‌شده در این فصل، یک ارزیابی جامع از روش پیشنهادی در سه بخش اصلی شامل ارزیابی طبقه‌بندی توییت‌ها، تولید سیگنال، و بک‌تست استراتژی را نشان می‌دهد. هر مرحله نقش حیاتی در ساخت یک چارچوب مقاوم ایفا می‌کند که بتواند احساسات شبکه‌های اجتماعی را به استراتژی‌های معاملاتی قابل اجرا تبدیل کند.

عملکرد مدل‌های طبقه‌بندی توییت‌ها مؤثر بودن اضافه کردن اطلاعات زمینه‌ای به مدل‌های زبانی را نشان می‌دهد. با گنجاندن هر دو زمینه‌ی بازار و زمانی در قالب پرسش‌ها، دقت طبقه‌بندی نسبت به روش‌های پایه به طور قابل توجهی بهبود یافته است. افزودن شاخص‌های بازاری نظیر نرخ بازدهی و شاخص قدرت روند درک دقیقی از پویایی‌های بازار ارائه می‌دهد، در حالی که اطلاعات زمانی مانع از تکیه‌ی بیش از حد مدل به روندهای قدیمی می‌شود. معیارهایی نظیر دقت، دقت، پوشش، و امتیاز \lr{F1}  به طور مداوم این بهبود را بازتاب داده و نشان‌دهنده قابلیت اطمینان خروجی‌های طبقه‌بندی هستند. علاوه بر این، تحلیل ماتریس درهم‌ریختگی کاهش قابل توجهی در خطاهای طبقه‌بندی در سه دسته‌ی احساسی صعودی، نزولی و خنثی را نشان می‌دهد. این مرحله به طور مؤثر پایه‌ای برای تولید سیگنال‌های دقیق معاملاتی ایجاد می‌کند و اطمینان حاصل می‌نماید که پیش‌بینی‌های احساسی با واقعیت‌های بازار تطابق نزدیکی دارند.

تبدیل احساسات سطح توییت به سیگنال‌های معاملاتی قابل اجرا نمایانگر مرحله‌ی حیاتی بعدی است. ارزیابی روش‌های تولید سیگنال، نقاط قوت دو رویکرد تجمیع شامل رأی‌گیری اکثریت و روش میانگین را نشان داد. در حالی که رأی‌گیری اکثریت سیگنال‌های پایدارتری با نویز کمتر ارائه داد، روش میانگین توانایی خود را در شناسایی تغییرات جزئی در احساسات تجمیع‌شده اثبات کرد. با استفاده از خروجی‌های طبقه‌بندی و ویژگی‌های بازار، سیگنال‌های تولیدشده توانستند پویایی‌های کوتاه‌مدت بازار را به طور مؤثری منعکس کنند. توانایی تولید سیگنال‌های متنوع اما قابل اجرا، انعطاف‌پذیری در برابر تغییرات رژیم‌های مختلف بازار را تضمین می‌کند و تعادلی بین دقت و غنای سیگنال ایجاد می‌نماید. این قابلیت در کاربردهای واقعی که نوسانات بازار و تغییرات احساسات نیاز به استراتژی‌های پویا دارند، بسیار مهم است.

در نهایت، نتایج بک‌تست استراتژی کاربرد عملی چارچوب پیشنهادی را در یک محیط معاملاتی واقعی برجسته می‌کند. در مقایسه با پرتفوی خرید و نگه‌داری، استراتژی پیشنهادی سازگاری و انعطاف‌پذیری برتری را در شرایط مختلف بازار نشان داد. در حالی که پرتفوی پایه به دلیل روند صعودی قوی بازار در ابتدا عملکرد بهتری داشت، در دوره‌های آشفتگی و افت ارزش با شکست مواجه شد. در مقابل، استراتژی پیشنهادی با بهره‌گیری از طیف گسترده‌ای از رژیم‌های بازار، بازده تجمعی برتری را ارائه داد. با مدیریت مؤثر ریسک‌ها در دوره‌های نزولی و بهره‌گیری از فرصت‌های کوتاه‌مدت، استراتژی توانایی خود را در حفظ رشد و افزایش سودآوری اثبات کرد. تحلیل جزئیات سود و زیان هر معامله نیز این مشاهده را تأیید نمود و نشان داد که استراتژی در سطح معامله موفق به حداکثرسازی سودها و حداقل‌سازی زیان‌ها شده است.

هم‌افزایی بین طبقه‌بندی دقیق توییت‌ها، تولید سیگنال‌های قوی و نتایج بک‌تست تطبیق‌پذیر، چارچوب جامعی برای معاملات مبتنی بر احساسات شبکه‌های اجتماعی ایجاد می‌کند. هر جزء دیگری را تکمیل می‌کند و یک زنجیره‌ی قدرتمند برای تبدیل احساسات به معاملات سودآور و قابل اجرا می‌سازد. نتایج نه تنها اثربخشی روش پیشنهادی را در مدیریت شرایط متنوع بازار تأیید می‌کنند، بلکه پتانسیل آن را برای کاربردهای واقعی در بازارهای مالی برجسته می‌سازند. یافته‌های ارائه‌شده در این فصل، پایه‌ای برای اکتشاف و بهبود بیشتر فراهم می‌کنند و نشان می‌دهند که ترکیب تکنیک‌های پیشرفته پردازش زبان طبیعی با مدل‌سازی مالی می‌تواند استراتژی‌های معاملاتی مبتنی بر احساسات را متحول کند.
